
%%FILE blueprint/src/chapters/algebra.tex
\chapter{The flag algebra $\mathcal{A}^\sigma$}
\label{chap:algebra}

Formal real combinations of flags form a free module; a product defined through
pair densities and a quotient by Razborov's averaging relations turn it into a
commutative $\mathbb{R}$-algebra, the flag algebra $\mathcal{A}^\sigma$.
(Source: \texttt{FlagAlgebra/FlagAlgebra.lean}.)

\begin{definition}[Size-tagged flag]
  \label{def:fin-flag}
  \lean{FlagAlgebras.FinFlag}
  \uses{def:flag-with-size}
  \leanok
  A \emph{size-tagged flag} is a dependent pair $\langle n, F\rangle$ of a vertex
  count $n$ and a flag $F$ on $\mathrm{Fin}\,n$. These are the basis indices of
  the algebra.
\end{definition}

\begin{definition}[Flag vector]
  \label{def:flag-vector}
  \lean{FlagAlgebras.FlagVector}
  \uses{def:fin-flag}
  \leanok
  A \emph{flag vector} is a finitely supported real combination of size-tagged
  flags, i.e.\ an element of $\mathrm{FinFlag}\,\sigma \to_{0} \mathbb{R}$: the
  underlying module before quotienting.
\end{definition}

\begin{definition}[Basis vector]
  \label{def:basis-vector}
  \lean{FlagAlgebras.basisVector}
  \uses{def:flag-vector}
  \leanok
  $\mathrm{basisVector}\,F$ is the unit vector supported on the single flag $F$;
  every construction reduces to these.
\end{definition}

\begin{definition}[Flag multiplication]
  \label{def:flag-mul}
  \lean{FlagAlgebras.flagMul, FlagAlgebras.flagMulWithSize}
  \uses{def:basis-vector, def:flag-density-n, def:flag-with-size}
  \leanok
  The product of two flags is expanded onto flags $G$ of a target size $\ell$,
  weighted by the pair density:
  \[ F \cdot F' \;=\; \sum_{G} d_2(F,F';G)\,\cdot G, \]
  taken at the minimal size $\ell = |F| + |F'| - |\sigma|$
  ($\mathrm{flagMul}$); the choice of $\ell$ is irrelevant modulo the averaging
  relations.
\end{definition}

\begin{definition}[Zero space]
  \label{def:zero-space}
  \lean{FlagAlgebras.ZeroSpace}
  \uses{def:zero-set, def:flag-vector}
  \leanok
  The \emph{zero space} $\mathrm{ZeroSpace}\,\sigma$ is the $\mathbb{R}$-submodule
  spanned by all of Razborov's averaging relations (the set
  $\mathrm{zeroSet}\,\sigma$). Quotienting by it identifies each flag with its
  density expansion.
\end{definition}

\begin{definition}[The flag algebra]
  \label{def:flag-algebra}
  \lean{FlagAlgebras.FlagAlgebra}
  \uses{def:flag-vector, def:zero-space, def:flag-vector-eqv}
  \leanok
  $\mathcal{A}^\sigma = \mathrm{FlagVector}\,\sigma / {\sim_v}$, flag vectors
  modulo the zero space. It carries a commutative $\mathbb{R}$-algebra structure
  (built below), with quotient map written $[\![\,\cdot\,]\!]$.
\end{definition}

\begin{theorem}[The zero space is an ideal]
  \label{thm:zero-space-ideal}
  \lean{FlagAlgebras.flagVector_mul_zeroSpace}
  \uses{def:zero-space, def:flag-mul}
  \leanok
  Multiplying any flag vector by an averaging relation gives another element of
  the zero space. Hence multiplication descends to a well-defined operation on
  the quotient.
\end{theorem}
\begin{proof}\leanok
  Reduce to a basis vector times a generator $\mathrm{zeroElement}\,F\,\ell$ and
  apply the chain rule (Theorem~\ref{thm:density-chain-rules}), which is exactly
  the statement that the product of the relation is again a telescoping sum of
  relations.
\end{proof}

\begin{theorem}[Associativity modulo $\sim_v$]
  \label{thm:mul-assoc}
  \lean{FlagAlgebras.flagVector_mul_assoc}
  \uses{def:flag-mul, def:flag-density-n}
  \leanok
  $(f\cdot g)\cdot h \sim_v f\cdot(g\cdot h)$. Associativity fails on the nose on
  flag vectors but holds modulo the zero space, so it lifts to genuine
  associativity on $\mathcal{A}^\sigma$.
\end{theorem}
\begin{proof}\leanok
  \uses{thm:triple-mul-density}
  Both sides expand, via the triple-density identity
  (Theorem~\ref{thm:triple-mul-density}), to the same sum
  $\sum_{G} d_3(F_1,F_2,F_3;G)\cdot G$.
\end{proof}

\begin{theorem}[Product equals density expansion]
  \label{thm:quot-mul-density}
  \lean{FlagAlgebras.basisVector_quot_mul_eq_flagMul_quot}
  \uses{def:flag-algebra, def:flag-mul}
  \leanok
  In $\mathcal{A}^\sigma$, $[\![F]\!]\cdot[\![G]\!] = [\![\,\mathrm{flagMul}\,F\,G\,]\!]$:
  the ring product of two flags is the class of their density expansion.
\end{theorem}
\begin{proof}\leanok
  Unfold the multiplication on the quotient, which is defined by
  $\mathrm{flagMul}$ and well-defined by Theorem~\ref{thm:zero-space-ideal}.
\end{proof}

\begin{theorem}[Same-size flags sum to one]
  \label{thm:sum-flags-eq-one}
  \lean{FlagAlgebras.sum_flagWithSize_eq_one}
  \uses{def:flag-algebra, def:basis-vector}
  \leanok
  For every $\ell \ge |\sigma|$, $\sum_{F : |F|=\ell} [\![F]\!] = 1$ in
  $\mathcal{A}^\sigma$: the sum over all flags of a fixed size is the unit. This
  normalization is used pervasively in density proofs.
\end{theorem}
\begin{proof}\leanok
  \uses{thm:basis-quot-eq-sum}
  Specialize the density-expansion identity
  (Theorem~\ref{thm:basis-quot-eq-sum}) to the empty flag, whose density in every
  size-$\ell$ flag is $1$.
\end{proof}

\section{Supporting constructions}

\begin{definition}[Flag of fixed size]
  \label{def:flag-with-size}
  \lean{FlagAlgebras.FlagWithSize}
  \uses{def:flag}
  \leanok
  $\mathrm{FlagWithSize}\,\sigma\,n = \mathrm{Flag}\,\sigma\,(\mathrm{Fin}\,n)$:
  a flag on the fixed $n$-element carrier.
\end{definition}

\begin{definition}[Density expansion]
  \label{def:flag-expansion}
  \lean{FlagAlgebras.flagExpansion}
  \uses{def:basis-vector, def:flag-density-n, def:flag-with-size}
  \leanok
  $\mathrm{flagExpansion}\,F\,\ell = \sum_{F'} d_1(F;F')\,\cdot F'$, the averaged
  expansion of $F$ onto size-$\ell$ flags.
\end{definition}

\begin{definition}[Averaging relation]
  \label{def:zero-element}
  \lean{FlagAlgebras.zeroElement}
  \uses{def:basis-vector, def:flag-expansion}
  \leanok
  Razborov's generating relation
  $\mathrm{zeroElement}\,F\,\ell = F - \mathrm{flagExpansion}\,F\,\ell$,
  identified with $0$ in the algebra.
\end{definition}

\begin{definition}[Relation set]
  \label{def:zero-set}
  \lean{FlagAlgebras.zeroSet}
  \uses{def:zero-element}
  \leanok
  The set of all averaging relations $\mathrm{zeroElement}\,F\,\ell$ (with
  $|F| \le \ell$); its span is the zero space.
\end{definition}

\begin{definition}[Flag equivalence of vectors]
  \label{def:flag-vector-eqv}
  \lean{FlagAlgebras.flagVectorEqv}
  \uses{def:zero-space}
  \leanok
  $f \sim_v g$ iff $f - g \in \mathrm{ZeroSpace}\,\sigma$: the two vectors
  represent the same flag-algebra element.
\end{definition}

\begin{theorem}[Quotient map is a homomorphism]
  \label{thm:quot-hom}
  \lean{FlagAlgebras.add_quot, FlagAlgebras.mul_quot, FlagAlgebras.smul_quot, FlagAlgebras.sum_quot}
  \uses{def:flag-vector, def:flag-algebra}
  \leanok
  The quotient map $[\![\,\cdot\,]\!]$ commutes with addition, multiplication,
  scalar multiplication, and finite sums.
\end{theorem}

\begin{theorem}[Triple product expansion]
  \label{thm:triple-mul-density}
  \lean{FlagAlgebras.three_flag_mul_eqv_sum_tripleDensity}
  \uses{def:basis-vector, def:flag-vector-eqv, def:flag-with-size, def:flag-density-n, def:flag-mul, thm:density-chain-rules}
  \leanok
  The product of three flags is, modulo $\sim_v$, the size-$\ell$ sum weighted by
  the triple density $d_3$. This is the engine behind associativity
  (Theorem~\ref{thm:mul-assoc}).
\end{theorem}

\begin{theorem}[Flag class as density expansion]
  \label{thm:basis-quot-eq-sum}
  \lean{FlagAlgebras.basisVector_quot_eq_sum}
  \uses{def:basis-vector, def:flag-with-size, def:flag-density-n, def:flag-algebra, thm:quot-hom, def:flag-expansion}
  \leanok
  In the algebra, $[\![F]\!] = \sum_{F'} d_1(F;F')\,[\![F']\!]$ over size-$\ell$
  flags: the averaging relation written as a $[\![\,\cdot\,]\!]$ identity.
\end{theorem}

\begin{theorem}[Nontriviality]
  \label{thm:flag-algebra-nontrivial}
  \lean{FlagAlgebras.flagAlgebra_zero_ne_one}
  \uses{def:flag-algebra, def:zero-space, def:basis-vector, def:flag-density-n, thm:density-chain-rules}
  \leanok
  $0 \ne 1$ in $\mathcal{A}^\sigma$: the flag algebra is nontrivial (so its
  $\mathrm{CommRing}$ structure is not the zero ring).
\end{theorem}
\begin{proof}\leanok
  If $1$ lay in $\mathrm{ZeroSpace}\,\sigma$ it would be a finite combination of
  averaging relations. But the linear functional $\varphi(G') = d_1(G';F)$
  (density into a fixed large flag $F$) sends every generator
  $\mathrm{zeroElement}\,G\,\ell$ to $0$ by the chain rule
  (Theorem~\ref{thm:density-chain-rules}), while $\varphi(1) = 1$ --- a
  contradiction.
\end{proof}


%%FILE blueprint/src/chapters/compute.tex
\chapter{The computational layer}
\label{chap:compute}

To actually run flag-algebra computations (enumerate flags, compute densities) at
elaboration time, the abstract objects are mirrored by executable ones: graphs
encoded as explicit edge sets $\mathrm{Finset}(\mathrm{Sym}_2(\mathrm{Fin}\,n))$.
Every computable definition is paired with a theorem proving it equals its
abstract counterpart, so computed results are provably correct. (Source:
\texttt{FlagAlgebra/Compute/}; all names in the \texttt{FlagAlgebras.Compute}
namespace.)

\begin{definition}[Edge-set graph]
  \label{def:sym2-graph}
  \lean{FlagAlgebras.Compute.Sym2Graph}
  \leanok
  A computable simple graph on $\mathrm{Fin}\,n$: a finite set of non-loop edges
  in $\mathrm{Sym}_2(\mathrm{Fin}\,n)$. It decodes to an abstract
  (empty-type) labeled graph.
\end{definition}

\begin{definition}[Edge-set labeled graph]
  \label{def:sym2-labeled-graph}
  \lean{FlagAlgebras.Compute.Sym2LabeledGraph}
  \uses{def:sym2-graph}
  \leanok
  The typed version: an edge set together with a graph embedding of a decoded
  type $\sigma$; decodes to an abstract $\mathrm{LabeledGraph}$.
\end{definition}

\begin{definition}[Computable flag]
  \label{def:sym2-flag}
  \lean{FlagAlgebras.Compute.Sym2Flag}
  \uses{def:sym2-labeled-graph, def:flag}
  \leanok
  The computable counterpart of $\mathrm{Flag}\,\sigma$: edge-set labeled graphs
  modulo a computable flag equivalence, carrying effective
  $\mathrm{Fintype}$/$\mathrm{DecidableEq}$ instances.
\end{definition}

\begin{definition}[Cheap iso invariant]
  \label{def:deg-key}
  \lean{FlagAlgebras.Compute.degKey}
  \uses{def:sym2-graph, def:deg-seq}
  \leanok
  A cheap $O(n)$ isomorphism invariant --- edge count paired with the sorted
  degree sequence --- used to prefilter before the expensive isomorphism test.
\end{definition}

\begin{theorem}[The key is an invariant]
  \label{thm:deg-key-iso-invariant}
  \lean{FlagAlgebras.Compute.degKey_iso_invariant}
  \uses{def:deg-key}
  \leanok
  Equivalent graphs share the cheap key, justifying its use as a deduplication
  prefilter.
\end{theorem}

\begin{definition}[Fast isomorphism test]
  \label{def:fast-iso}
  \lean{FlagAlgebras.Compute.isEmptyIsoFast_bool, FlagAlgebras.Compute.isIsoFast_bool}
  \uses{def:sym2-labeled-graph, def:build-full-map}
  \leanok
  Permutation-search isomorphism decision procedures for edge-set graphs (empty
  and $\sigma$-typed), early-exiting on unequal keys; each is proved sound and
  complete against flag equivalence.
\end{definition}

\begin{definition}[Fast decidability instance]
  \label{def:fast-iso-decidable}
  \lean{FlagAlgebras.Compute.fastDecidableSym2LabeledGraphEqv}
  \uses{def:fast-iso}
  \leanok
  The high-priority $\mathrm{Decidable}$ instance for flag equivalence backed by
  the fast test; it makes $\mathrm{Fintype}$-based flag enumeration tractable.
\end{definition}

\begin{theorem}[Computable normalizing factor, certified]
  \label{thm:downward-compute}
  \lean{FlagAlgebras.Compute.downwardNormalizingFactor_Sym2Flag, FlagAlgebras.Compute.downwardNormalizingFactor_eq}
  \uses{def:sym2-flag, def:downward-factor}
  \leanok
  A computable downward normalizing factor on edge-set flags, proved equal to the
  abstract $\mathrm{downwardNormalizingFactor}$ on the decoded flag.
\end{theorem}

\begin{theorem}[Computable densities, certified]
  \label{thm:density-compute}
  \lean{FlagAlgebras.Compute.flagDensity₁_eq_sym2EmptyTypeFlagDensity₁, FlagAlgebras.Compute.flagDensity₂_eq_sym2FlagDensity₂}
  \uses{def:sym2-flag, def:flag-density-n, def:sym2-density1, def:sym2-density2}
  \leanok
  Computable single- and pair-densities on edge-set flags, each proved equal to
  the abstract $d_1$, $d_2$ --- so a density computed in the edge-set
  representation is provably the abstract value.
\end{theorem}

\begin{theorem}[Flag enumeration, complete]
  \label{thm:flag-enumeration}
  \lean{FlagAlgebras.Compute.genSym2Graphs, FlagAlgebras.Compute.genFlagSet_eq_univ}
  \uses{def:sym2-graph, def:fast-iso, def:deg-key}
  \leanok
  Orderly (McKay-style) augmentation with cheap-key dedup generates one
  representative per isomorphism class; the generated set is proved to exhaust
  $\mathrm{Finset.univ}$, so enumeration is complete.
\end{theorem}

\section{Supporting constructions}

\begin{definition}[Computable empty-typed flag]
  \label{def:sym2-empty-flag}
  \lean{FlagAlgebras.Compute.Sym2EmptyTypedFlag}
  \uses{def:sym2-graph}
  \leanok
  The computable counterpart of $\mathrm{Flag}\,\emptyset\,(\mathrm{Fin}\,n)$:
  edge-set graphs modulo computable flag equivalence.
\end{definition}

\begin{definition}[Placement counts]
  \label{def:sym2-count}
  \lean{FlagAlgebras.Compute.sym2InducedSubgraphListCount, FlagAlgebras.Compute.sym2InducedLabeledSubgraphListCount}
  \uses{def:sym2-graph, def:sym2-labeled-graph}
  \leanok
  The computable numerators of the densities: the number of vertex-disjoint
  placements of a flag list inside a host, empty-typed and $\sigma$-typed.
\end{definition}

\begin{definition}[Computable single density]
  \label{def:sym2-density1}
  \lean{FlagAlgebras.Compute.sym2EmptyTypeFlagDensity₁}
  \uses{def:sym2-empty-flag, def:sym2-count}
  \leanok
  The computable single subflag density on empty-typed edge-set flags.
\end{definition}

\begin{definition}[Computable pair density]
  \label{def:sym2-density2}
  \lean{FlagAlgebras.Compute.sym2FlagDensity₂}
  \uses{def:sym2-flag, def:sym2-count}
  \leanok
  The computable pair (multiplication) density on $\sigma$-typed edge-set flags.
\end{definition}

\begin{definition}[Degree sequence]
  \label{def:deg-seq}
  \lean{FlagAlgebras.Compute.degSeq}
  \uses{def:sym2-graph}
  \leanok
  The sorted degree sequence of an edge-set graph, the main part of the cheap
  isomorphism key.
\end{definition}

\begin{definition}[Permutation builder]
  \label{def:build-full-map}
  \lean{FlagAlgebras.Compute.buildFullMap}
  \leanok
  Assembles a full vertex permutation that fixes the type embedding and routes the
  remaining vertices through a chosen permutation --- the inner loop of the fast
  isomorphism search.
\end{definition}


%%FILE blueprint/src/chapters/density.tex
\chapter{Subflag densities}
\label{chap:density}

The density of one flag inside another is the typed analogue of induced-subgraph
density. This chapter defines it, lifts it through the flag quotient, generalizes
it to lists of disjoint flags, and proves the \emph{chain rules} that make the
whole flag algebra tick. (Source: \texttt{FlagAlgebra/SubflagDensity.lean},
\texttt{SubflagListDensity.lean}, \texttt{SubflagListDensityProp.lean}.)

\begin{definition}[Labeled-graph count]
  \label{def:labeled-graph-count}
  \lean{FlagAlgebras.labeledGraphCount}
  \uses{def:labeled-graph}
  \leanok
  $\mathrm{labeledGraphCount}\,H\,G$ is the number of induced labeled subgraphs
  of the host $G$ that are labeled-isomorphic to $H$.
\end{definition}

\begin{definition}[Labeled-graph density]
  \label{def:labeled-graph-density}
  \lean{FlagAlgebras.labeledGraphDensity}
  \uses{def:labeled-graph-count}
  \leanok
  $\mathrm{labeledGraphDensity}\,H\,G$ normalizes the count by the number of ways
  to choose the non-type vertices of an $H$-sized induced subgraph, i.e.\ by
  $\binom{|G|-|\sigma|}{|H|-|\sigma|}$, giving a rational in $[0,1]$.
\end{definition}

\begin{lemma}[Isomorphism invariance]
  \label{lem:density-iso-invariant}
  \lean{FlagAlgebras.labeledGraphDensity_respect_eqv}
  \uses{def:labeled-graph-density, def:labeled-graph-iso}
  \leanok
  Labeled-graph density is invariant under $\simeq_f$ in both arguments. This is
  the well-definedness fact that lets density descend to flags.
\end{lemma}
\begin{proof}\leanok
  Transport induced copies along the given isomorphisms; the counts, and hence
  the normalized densities, agree.
\end{proof}

\begin{definition}[Subflag density]
  \label{def:subflag-density}
  \lean{FlagAlgebras.subflagDensity}
  \uses{def:labeled-graph-density, lem:density-iso-invariant, def:flag}
  \leanok
  $\mathrm{subflagDensity} : \mathrm{Flag}\,\sigma\,V \to \mathrm{Flag}\,\sigma\,W
  \to \mathbb{Q}$ is $\mathrm{labeledGraphDensity}$ lifted through the flag
  quotient in both arguments: the density of one flag inside another.
\end{definition}

\begin{lemma}[Boundary values]
  \label{lem:subflag-density-values}
  \lean{FlagAlgebras.subflagDensity_empty, FlagAlgebras.subflagDensity_self, FlagAlgebras.subflagDensity_other}
  \uses{def:subflag-density, def:empty-flag}
  \leanok
  The empty flag has density $1$ in every flag; a flag has density $1$ in
  itself; two distinct flags of the same size have density $0$ in each other.
\end{lemma}

\begin{definition}[List density and fixed arities]
  \label{def:flag-list-density}
  \lean{FlagAlgebras.flagListDensity}
  \uses{def:flag-list}
  \leanok
  $\mathrm{flagListDensity}$ counts the induced subgraph-lists of $G$ that
  simultaneously realize a whole disjoint list of prescribed flags (normalized by
  the multinomial coefficient distributing the non-type vertices), lifted through
  the quotient. It generalizes $\mathrm{labeledGraphListDensity}$ on carriers.
\end{definition}

\begin{definition}[Single, pair, and triple densities]
  \label{def:flag-density-n}
  \lean{FlagAlgebras.flagDensity₁, FlagAlgebras.flagDensity₂, FlagAlgebras.flagDensity₃}
  \uses{def:flag-list-density}
  \leanok
  The fixed-arity specializations
  $d_1(F;G)$, $d_2(F_1,F_2;G)$, $d_3(F_1,F_2,F_3;G)$ are list densities of the
  singleton, pair, and triple lists. $d_2$ is the factor multiplication is built
  from; $d_1$ is the ordinary subflag density.
\end{definition}

\begin{theorem}[Single density is subflag density]
  \label{thm:subflag-eq-list}
  \lean{FlagAlgebras.subflagDensity_eq_flagListDensity}
  \uses{def:subflag-density, def:flag-density-n}
  \leanok
  $\mathrm{subflagDensity}\,F\,G = d_1(F;G)$: the single-flag list density
  coincides with the subflag density of Chapter start, linking the two
  developments.
\end{theorem}
\begin{proof}\leanok
  Both count the same induced copies; unfold the singleton list.
\end{proof}

\begin{lemma}[Permutation and empty-insertion invariance]
  \label{lem:density-invariance}
  \lean{FlagAlgebras.flagDensity_permute, FlagAlgebras.flagDensity_insert_empty}
  \uses{def:flag-list-density}
  \leanok
  List density is unchanged by permuting the list, and by inserting an empty
  flag. These symmetries reduce the fixed-arity densities to one another in the
  chain-rule proofs.
\end{lemma}

\begin{theorem}[Density chain rules]
  \label{thm:density-chain-rules}
  \lean{FlagAlgebras.flagDensity_eq_sum_density_prods, FlagAlgebras.flagPairDensity_eq_sum_density_prods, FlagAlgebras.flagPairDensity_eq_sum_density_prods', FlagAlgebras.flagTripleDensity_eq_sum_density_prods}
  \uses{def:flag-density-n, lem:density-invariance}
  \leanok
  A product-flag density expands as a sum over intermediate flags $G'$ of size
  $\ell'$ of products of densities, e.g.
  \[ d_1(F;G) = \sum_{G'} d_1(F;G')\, d_1(G';G), \qquad
     d_2(F_1,F_2;G) = \sum_{G'} d_1(F_1;G')\, d_2(G',F_2;G). \]
  The triple rule is foundational; the pair and single rules descend from it.
  These are the identities the quotient algebra is defined to respect.
\end{theorem}
\begin{proof}\leanok
  Partition the induced copies of the product by the intermediate flag they span
  on a size-$\ell'$ vertex set, then count each part with the appropriate
  density.
\end{proof}

\begin{theorem}[Asymptotic independence]
  \label{thm:density-prod-approx}
  \lean{FlagAlgebras.flagListDensity₂_prod_approx}
  \uses{def:flag-density-n, thm:subflag-eq-list}
  \leanok
  There is a constant $c = 2(|F|+|F'|)^2$, independent of $G$, with
  \[ \bigl| d_2(F,F';G) - d_1(F;G)\, d_1(F';G) \bigr| \le \frac{c}{|G|}. \]
  Distinct subflag densities are thus asymptotically multiplicative --- the
  probabilistic heart of the ``product of densities'' reasoning.
\end{theorem}
\begin{proof}\leanok
  Model the choice of two vertex subsets as a finite probability space; the
  discrepancy is the probability that the subsets overlap off the type, which is
  $O(1/|G|)$ by a union bound.
\end{proof}


%%FILE blueprint/src/chapters/downward.tex
\chapter{The unlabeling operator}
\label{chap:downward}

Density bounds are ultimately statements about \emph{unlabeled} graphs. The
\emph{downward} operator $[\![\,\cdot\,]\!]_0 : \mathcal{A}^\sigma \to
\mathcal{A}^\emptyset$ forgets a flag's labelling by averaging over all ways to
place the labels. (Source: \texttt{FlagAlgebra/FlagOperators.lean}; $\emptyset$
denotes the empty type $\varnothing_t$.)

\begin{definition}[Unlabel]
  \label{def:unlabel}
  \lean{FlagAlgebras.unlabel}
  \uses{def:flag, def:unlabeled-graph, def:empty-type}
  \leanok
  $\mathrm{unlabel}$ drops the type embedding of a single flag, returning its
  underlying graph as an unlabeled ($\emptyset$-type) flag.
\end{definition}

\begin{definition}[Downward normalizing factor]
  \label{def:downward-factor}
  \lean{FlagAlgebras.downwardNormalizingFactor}
  \uses{def:flag}
  \leanok
  The strictly positive rational weight of a flag: the fraction of label
  placements (injections of the $|\sigma|$ labels into the vertices) that realize
  it. It is the coefficient by which $\mathrm{unlabel}\,F$ is scaled.
\end{definition}

\begin{lemma}[The bridge set, counted three ways]
  \label{lem:bridge-counts}
  \lean{FlagAlgebras.isoInjectiveMapSet_card_eq_isomorphismCount_mul_labeledGraphCount, FlagAlgebras.isoInjectiveMapSet_card_eq_sum_labeledGraphCount_of_same_graph, FlagAlgebras.isoInjectiveMapSet_card_eq_sum_labelExtensions_isomorphismCount_mul_labeledGraphCount}
  \uses{def:unlabel, def:isomorphism-count}
  \leanok
  The set of pairs (vertex subset, injective label assignment) embedding a
  labeled flag into an unlabeled one is counted three ways: by
  isomorphism-count times labeled count; by summing labeled counts over labellings
  of the same graph; and, grouped by label-extension flag, over all relabellings.
\end{lemma}

\begin{theorem}[Counting identity]
  \label{thm:counting-identity}
  \lean{FlagAlgebras.flagDensity_mul_downwardNormalizingFactor_eq_sum_labelExtensions}
  \uses{lem:bridge-counts, def:unlabel, def:downward-factor, def:flag-density-n, def:label-extensions}
  \leanok
  Equating the three counts yields
  \[ d_1(\mathrm{unlabel}\,F;\,F')\cdot w(F)
     \;=\; \sum_{G \in \mathrm{ext}(F')} d_1(F;G)\cdot w(G), \]
  where $w$ is the normalizing factor and $\mathrm{ext}(F')$ ranges over the
  relabellings of $F'$. This is the algebraic engine of well-definedness.
\end{theorem}
\begin{proof}\leanok
  Combine the first and third counts of Lemma~\ref{lem:bridge-counts} and divide
  by the number of label injections.
\end{proof}

\begin{lemma}[Downward respects the zero space]
  \label{lem:downward-well-defined}
  \lean{FlagAlgebras.downwardFlagVector_zeroSpace}
  \uses{def:zero-space}
  \leanok
  The unlabeling map sends the zero space of $\sigma$ into the zero space of the
  empty type, so it descends to the quotient algebra.
\end{lemma}
\begin{proof}\leanok
  \uses{lem:downward-zero-element}
  The zero space is spanned by the generating relations, each of which downward
  sends into the zero space of the empty type
  (Lemma~\ref{lem:downward-zero-element}).
\end{proof}

\begin{definition}[Downward operator]
  \label{def:downward}
  \lean{FlagAlgebras.downward}
  \uses{def:unlabel, def:downward-factor, lem:downward-well-defined, def:flag-algebra, def:downward-flag}
  \leanok
  $[\![\,\cdot\,]\!]_0 : \mathcal{A}^\sigma \to \mathcal{A}^\emptyset$ is the
  $\mathbb{R}$-linear map induced by $F \mapsto w(F)\cdot \mathrm{unlabel}\,F$;
  it averages a labeled element over all unlabelings.
\end{definition}

\begin{theorem}[Linearity]
  \label{thm:downward-linear}
  \lean{FlagAlgebras.downward_add, FlagAlgebras.downward_smul}
  \uses{def:downward}
  \leanok
  $[\![f+f']\!]_0 = [\![f]\!]_0 + [\![f']\!]_0$ and
  $[\![r\cdot f]\!]_0 = r\cdot[\![f]\!]_0$: downward is $\mathbb{R}$-linear.
\end{theorem}

\begin{theorem}[Identity on the empty type]
  \label{thm:downward-empty}
  \lean{FlagAlgebras.downward_emptyType}
  \uses{def:downward}
  \leanok
  On $\mathcal{A}^\emptyset$ there is nothing to unlabel, so
  $[\![f]\!]_0 = f$.
\end{theorem}

\section{Supporting constructions}

\begin{definition}[Empty type]
  \label{def:empty-type}
  \lean{FlagAlgebras.emptyType}
  \uses{def:flag-type}
  \leanok
  The empty flag type $\varnothing_t$ (written $\emptyset$) with no labeled
  vertices; its flags are unlabeled graphs, and density bounds live in
  $\mathcal{A}^\emptyset$.
\end{definition}

\begin{definition}[Isomorphism count]
  \label{def:isomorphism-count}
  \lean{FlagAlgebras.isomorphismCount}
  \uses{def:labeled-graph}
  \leanok
  The number of label placements on the same underlying graph as $G$ that are
  flag-isomorphic to $G$.
\end{definition}

\begin{definition}[Unlabeled graph]
  \label{def:unlabeled-graph}
  \lean{FlagAlgebras.unlabeledGraph, FlagAlgebras.unlabeledGraphQuot}
  \uses{def:labeled-graph, def:empty-type}
  \leanok
  Forgetting the type embedding of a labeled graph, as a graph
  ($\mathrm{unlabeledGraph}$) and as its flag class
  ($\mathrm{unlabeledGraphQuot}$).
\end{definition}

\begin{definition}[Label extensions]
  \label{def:label-extensions}
  \lean{FlagAlgebras.labelExtensions}
  \uses{def:unlabel, def:flag-with-size, def:empty-type}
  \leanok
  For an unlabeled flag $F$, the finset of all $\sigma$-labeled flags of the same
  size whose unlabeling is $F$ --- the ways to re-add a labelling.
\end{definition}

\begin{definition}[Downward image of a flag]
  \label{def:downward-flag}
  \lean{FlagAlgebras.downwardFlag, FlagAlgebras.downwardFlagVector, FlagAlgebras.downwardFlagVectorQuot}
  \uses{def:downward-factor, def:unlabel, def:basis-vector}
  \leanok
  $\mathrm{downwardFlag}\,F = w(F)\cdot\mathrm{unlabel}\,F$ on a single flag, its
  linear extension $\mathrm{downwardFlagVector}$ to flag vectors, and the
  quotient-level $\mathrm{downwardFlagVectorQuot}$ --- the raw form of downward
  before descending to the quotient.
\end{definition}

\begin{lemma}[Downward averaging expansion]
  \label{lem:downward-flag-sum}
  \lean{FlagAlgebras.downwardFlag_eqv_sum_flagDensity_smul_downwardFlag}
  \uses{def:downward-flag, thm:counting-identity, def:flag-density-n, def:label-extensions}
  \leanok
  $\mathrm{downwardFlag}\,F \sim_v \sum_{G} d_1(F;G)\cdot\mathrm{downwardFlag}\,G$
  over size-$\ell$ flags --- the averaging expansion driving well-definedness.
\end{lemma}

\begin{lemma}[Downward of a relation is a relation]
  \label{lem:downward-zero-element}
  \lean{FlagAlgebras.downwardFlagVector_zeroElement_zeroSpace}
  \uses{def:downward-flag, lem:downward-flag-sum, def:zero-space, def:zero-element}
  \leanok
  $\mathrm{downwardFlagVector}$ sends each generating relation
  $\mathrm{zeroElement}\,F\,\ell$ into the zero space of the empty type.
\end{lemma}

\begin{theorem}[The normalizing factor is positive]
  \label{lem:downward-factor-pos}
  \lean{FlagAlgebras.downwardNormalizingFactor_pos}
  \uses{def:downward-factor, def:isomorphism-count}
  \leanok
  $w(F) > 0$ for every flag $F$: every flag is realized by at least one label
  placement.
\end{theorem}


%%FILE blueprint/src/chapters/flags.tex
\chapter{Flags and labeled graphs}
\label{chap:flags}

Fix a finite label set $T$ and a graph $\sigma$ on it, the \emph{type}. A flag of
type $\sigma$ is a finite graph carrying a distinguished induced copy of $\sigma$,
taken up to type-preserving isomorphism. This chapter sets up these carriers;
everything downstream is built from them. (Source: \texttt{FlagAlgebra/FlagDef.lean}.)

\begin{definition}[Flag type]
  \label{def:flag-type}
  \lean{FlagAlgebras.FlagType}
  \leanok
  A \emph{flag type} $\sigma$ is a simple graph on a finite label set $T$. It
  describes the labeled ``root'' every flag of type $\sigma$ must carry.
\end{definition}

\begin{definition}[Labeled graph]
  \label{def:labeled-graph}
  \lean{FlagAlgebras.LabeledGraph}
  \uses{def:flag-type}
  \leanok
  A \emph{labeled graph} of type $\sigma$ on a vertex type $V$ is a simple graph
  $G$ on $V$ together with an induced-subgraph embedding
  $\sigma \hookrightarrow G$. Its image is the set of labeled (``rooted'')
  vertices.
\end{definition}

\begin{definition}[Labeled subgraph]
  \label{def:labeled-subgraph}
  \lean{FlagAlgebras.LabeledSubgraph}
  \uses{def:labeled-graph}
  \leanok
  A \emph{labeled subgraph} of a labeled graph $G$ is a subgraph of $G$ that
  still contains all labeled vertices and embeds $\sigma$ compatibly with $G$.
  These are the pieces cut out when measuring subflag densities.
\end{definition}

\begin{definition}[Isomorphism of labeled graphs, $\simeq_f$]
  \label{def:labeled-graph-iso}
  \lean{FlagAlgebras.LabeledGraphIso}
  \uses{def:labeled-graph}
  \leanok
  A \emph{labeled-graph isomorphism} $G \simeq_f G'$ is a graph isomorphism of the
  underlying graphs that commutes with the two type embeddings. Flags are exactly
  the isomorphism classes for this relation.
\end{definition}

\begin{definition}[Flag]
  \label{def:flag}
  \lean{FlagAlgebras.Flag}
  \uses{def:labeled-graph, def:labeled-graph-iso, def:flag-eqv}
  \leanok
  A \emph{flag} of type $\sigma$ on $V$, $\mathrm{Flag}\,\sigma\,V$, is an
  isomorphism class of labeled graphs under $\sim_f$ (the equivalence
  ``$\exists$ a labeled-graph isomorphism''). This is the basic object the flag
  algebra is built from.
\end{definition}

\begin{definition}[Empty flag]
  \label{def:empty-flag}
  \lean{FlagAlgebras.emptyFlag}
  \uses{def:flag, def:empty-labeled-graph}
  \leanok
  The \emph{empty flag} $\mathrm{emptyFlag}\,\sigma$ is the class of the trivial
  labeled graph $\sigma$ on $T$ with the identity embedding: the type with no
  extra vertices. It is the multiplicative unit of the flag algebra.
\end{definition}

\begin{definition}[Flag list]
  \label{def:flag-list}
  \lean{FlagAlgebras.FlagList}
  \uses{def:flag}
  \leanok
  A \emph{flag list} of length $t$ is a tuple $(F_0,\dots,F_{t-1})$ of flags of
  type $\sigma$, the $i$-th on its own vertex type. Flag lists (with the
  type-level \texttt{insert}/\texttt{permute} operations) are what let us speak
  of the simultaneous density of several disjoint flags.
\end{definition}

\section{Supporting constructions}

\begin{definition}[Empty labeled graph]
  \label{def:empty-labeled-graph}
  \lean{FlagAlgebras.emptyLabeledGraph}
  \uses{def:labeled-graph}
  \leanok
  The trivial carrier: the type $\sigma$ itself on $T$ with the identity
  embedding. Its isomorphism class is the empty flag.
\end{definition}

\begin{definition}[Flag equivalence]
  \label{def:flag-eqv}
  \lean{FlagAlgebras.flagEqv}
  \uses{def:labeled-graph, def:labeled-graph-iso}
  \leanok
  The relation $G \sim_f G'$: ``there exists a labeled-graph isomorphism
  $G \simeq_f G'$''. Flags are its equivalence classes.
\end{definition}

\begin{definition}[Induced labeled subgraph]
  \label{def:induced-labeled-subgraph}
  \lean{FlagAlgebras.LabeledSubgraph.inducedLabeledSubgraph}
  \uses{def:labeled-graph, def:labeled-subgraph}
  \leanok
  The labeled subgraph of $G$ induced on a chosen vertex set $S$ containing all
  labeled vertices --- the canonical way to cut out a sub-flag on $S$.
\end{definition}

\begin{definition}[Labeled-graph list]
  \label{def:labeled-graph-list}
  \lean{FlagAlgebras.LabeledGraphList}
  \uses{def:labeled-graph}
  \leanok
  A $t$-tuple of labeled graphs, the $i$-th on vertex type $V_i$; the
  pre-quotient form of a flag list.
\end{definition}

\begin{definition}[Type-level list operations]
  \label{def:list-type-ops}
  \lean{FlagAlgebras.listTypeInsert, FlagAlgebras.listTypePermute}
  \leanok
  The operations on a vertex-type family that append a new type at the end
  ($\mathrm{listTypeInsert}$) or reorder it by a permutation
  ($\mathrm{listTypePermute}$); they supply the return types of the flag-list
  operations.
\end{definition}

\begin{definition}[Flag-list operations]
  \label{def:flag-list-ops}
  \lean{FlagAlgebras.FlagList.insert, FlagAlgebras.FlagList.permute}
  \uses{def:flag-list, def:list-type-ops}
  \leanok
  Appending a flag to a flag list, and reordering a flag list by a permutation.
\end{definition}

\begin{theorem}[Quotient of lists is a list of flags]
  \label{thm:eqv-quot-list-flaglist}
  \lean{FlagAlgebras.eqv_QuotLabeledGraphList_FlagList}
  \uses{def:labeled-graph-list, def:flag-list, def:flag-eqv}
  \leanok
  The quotient of labeled-graph lists is equivalent to the tuple of flags: taking
  isomorphism classes commutes with the $\mathrm{Fin}\,t$-indexed product.
\end{theorem}


%%FILE blueprint/src/chapters/limits.tex
\chapter{Limit semantics and random homomorphisms}
\label{chap:limits}

The analytic half of the theory identifies the algebraic positive homomorphisms
with genuine limits of subgraph densities (Razborov's Theorem~3.3), then builds a
probability measure on them (Theorem~3.5) that yields the Cauchy--Schwarz
inequality powering SDP-style bounds. (Source:
\texttt{FlagAlgebra/FlagSequence.lean}, \texttt{RandomHom.lean}.)

\begin{definition}[Flag sequence]
  \label{def:flag-seq}
  \lean{FlagAlgebras.FlagSeq}
  \uses{def:fin-flag}
  \leanok
  A \emph{flag sequence} is a sequence $\mathbb{N} \to \mathrm{FinFlag}\,\sigma$
  of flags, typically of growing size.
\end{definition}

\begin{definition}[Density space]
  \label{def:flag-density-space}
  \lean{FlagAlgebras.FlagDensitySpace}
  \uses{def:fin-flag}
  \leanok
  The compact, metrizable space of density profiles $\mathrm{FinFlag}\,\sigma \to
  [0,1]$ (a Tychonoff cube); the ambient space of the limit semantics.
\end{definition}

\begin{definition}[Convergence]
  \label{def:converges-to}
  \lean{FlagAlgebras.ConvergesTo}
  \uses{def:flag-seq, def:flag-density-space, def:increases, def:flag-density-seq}
  \leanok
  $\mathrm{ConvergesTo}\,s\,a$ says the flag sequence $s$ has strictly increasing
  sizes and its subgraph-density profile converges pointwise to the limit $a$.
\end{definition}

\begin{definition}[Density profile of a homomorphism]
  \label{def:positive-hom-coe}
  \lean{FlagAlgebras.PositiveHom.coe}
  \uses{def:positive-hom, def:flag-density-space}
  \leanok
  Each positive homomorphism $\varphi$ has a density profile
  $F \mapsto \varphi[\![F]\!]$; this is the injective bridge from the algebraic
  to the analytic side.
\end{definition}

\begin{definition}[Positive homomorphism space]
  \label{def:positive-hom-space}
  \lean{FlagAlgebras.PositiveHomSpace}
  \uses{def:positive-hom-coe}
  \leanok
  The range of the profile map; shown to be closed, hence compact and
  measurable.
\end{definition}

\begin{theorem}[Algebraic characterization]
  \label{thm:positive-hom-space-char}
  \lean{FlagAlgebras.positiveHomSpace_eq, FlagAlgebras.zeroSpaceProp, FlagAlgebras.oneProp, FlagAlgebras.mulProp}
  \uses{def:positive-hom-space, def:flag-density-n}
  \leanok
  A density profile arises from a positive homomorphism iff it satisfies the
  chain rule (\texttt{zeroSpaceProp}), normalization $a(1)=1$ (\texttt{oneProp}),
  and multiplicativity through pair densities (\texttt{mulProp}).
\end{theorem}

\begin{theorem}[Every limit is a homomorphism (Razborov 3.3a)]
  \label{thm:limit-is-hom}
  \lean{FlagAlgebras.flagSeq_limit_mem_positiveHom}
  \uses{def:converges-to, def:positive-hom-coe, thm:positive-hom-space-char}
  \leanok
  Any convergent flag-sequence density limit is the profile of some positive
  homomorphism.
\end{theorem}
\begin{proof}\leanok
  The limit inherits the three structural properties of
  Theorem~\ref{thm:positive-hom-space-char} from the finite densities, so it lies
  in the positive-homomorphism space.
\end{proof}

\begin{theorem}[Every homomorphism is a limit (Razborov 3.3b)]
  \label{thm:hom-is-limit}
  \lean{FlagAlgebras.positiveHom_as_flagSeq_limit}
  \uses{def:converges-to, def:positive-hom-coe, def:flagseq-measure}
  \leanok
  Conversely every positive homomorphism arises as a convergent flag-sequence
  limit.
\end{theorem}
\begin{proof}\leanok
  Sample random flags from the profile and apply a Chebyshev / Borel--Cantelli
  second-moment argument to get almost-sure convergence to $\varphi$.
\end{proof}

\begin{definition}[Random flag distribution]
  \label{def:finflag-pmf}
  \lean{FlagAlgebras.FinFlag.toPMF}
  \uses{def:flag-density-space}
  \leanok
  Turns a large unlabeled flag of positive $\sigma$-density into a distribution
  on density profiles: draw a random labelling extension weighted by its
  normalizing factor and read off its density vector.
\end{definition}

\begin{theorem}[Random homomorphism measure (Razborov 3.5)]
  \label{thm:random-hom-measure}
  \lean{FlagAlgebras.exists_probMeasure_extend_emptyType_positiveHom}
  \uses{def:finflag-pmf, def:positive-hom-space, thm:positive-hom-space-char, thm:hom-is-limit, def:downward, thm:lemma-3-11, lem:measure-support}
  \leanok
  For a base homomorphism $\varphi_0$ on the empty type with
  $\varphi_0\langle\sigma\rangle_0 > 0$, there is a probability measure
  $\mathbb{P}[\varphi_0]$ on the positive-homomorphism space whose expectation
  recovers the normalized downward operator:
  \[ \int \varphi(f)\, d\mathbb{P}[\varphi_0]
     \;=\; \frac{\varphi_0[\![f]\!]_0}{\varphi_0[\![1]\!]_0}. \]
\end{theorem}
\begin{proof}\leanok
  Build $\mathbb{P}[\varphi_0]$ as a Prokhorov weak limit of the finite
  distributions of Definition~\ref{def:finflag-pmf} along a converging flag
  sequence; its support lies in the positive-homomorphism space, and the integral
  is computed on the dense finite level.
\end{proof}

\begin{theorem}[Downward preserves the semantic cone]
  \label{thm:downward-preserve-cone}
  \lean{FlagAlgebras.downward_preserve_semanticCone}
  \uses{thm:random-hom-measure, def:semantic-cone, def:downward}
  \leanok
  If $f$ is semantically nonnegative in the typed algebra then so is its
  unlabeling in the empty-type algebra:
  \[ f \in \mathrm{semanticCone}\,\sigma \;\Longrightarrow\;
     [\![f]\!]_0 \in \mathrm{semanticCone}\,\emptyset. \]
  This is the bridge that turns typed nonnegativity into unlabeled
  nonnegativity, and it is what the sum-of-squares certificates
  (Theorem~\ref{thm:sos-downward-nonneg}) and the square/Cauchy--Schwarz results
  below rest on.
\end{theorem}
\begin{proof}\leanok
  \uses{lem:downward-zero-at-hom, lem:positivehom-one-downward-pos, thm:prob-measure-spec}
  Fix a base homomorphism $\varphi_0$. If $\varphi_0\langle\sigma\rangle_0 = 0$
  then $\varphi_0[\![f]\!]_0 = 0$ (Lemma~\ref{lem:downward-zero-at-hom}).
  Otherwise it is positive, and Theorem~\ref{thm:prob-measure-spec} expresses
  $\varphi_0[\![f]\!]_0$ as the positive multiple
  $\varphi_0[\![1]\!]_0 \cdot \int \varphi(f)\, d\mathbb{P}[\varphi_0]$
  (Lemma~\ref{lem:positivehom-one-downward-pos}) of the integral of a nonnegative
  integrand, hence $\ge 0$.
\end{proof}

\begin{theorem}[Squares have nonnegative average]
  \label{thm:square-downward-nonneg}
  \lean{FlagAlgebras.square_downward_nonneg}
  \uses{def:downward}
  \leanok
  $[\![f\cdot f]\!]_0 \ge 0$ for every $f$: the downward image of a square is
  nonnegative.
\end{theorem}
\begin{proof}\leanok
  \uses{thm:downward-preserve-cone}
  Every homomorphism sends $f\cdot f$ to $\varphi(f)^2 \ge 0$, so $f\cdot f$ is in
  the semantic cone; apply Theorem~\ref{thm:downward-preserve-cone}.
\end{proof}

\begin{theorem}[Cauchy--Schwarz for downward]
  \label{thm:cauchy-schwarz}
  \lean{FlagAlgebras.square_downward_mul_ge_mul_downward_square}
  \uses{thm:random-hom-measure, def:positive-hom-space}
  \leanok
  $[\![f\cdot f]\!]_0 \cdot [\![g\cdot g]\!]_0 \ge [\![f\cdot g]\!]_0^2$, the key
  inequality behind semidefinite flag-algebra density bounds.
\end{theorem}
\begin{proof}\leanok
  \uses{thm:prob-measure-spec, lem:downward-zero-at-hom}
  Apply the $L^2$ Cauchy--Schwarz (Hölder) inequality to the evaluation functions
  $\varphi \mapsto \varphi(f), \varphi(g)$ under $\mathbb{P}[\varphi_0]$
  (Theorem~\ref{thm:prob-measure-spec}), handling the degenerate base case with
  Lemma~\ref{lem:downward-zero-at-hom}.
\end{proof}

\section{Supporting constructions}

\begin{definition}[Increasing sizes]
  \label{def:increases}
  \lean{FlagAlgebras.Increases}
  \uses{def:flag-seq}
  \leanok
  A flag sequence has strictly increasing flag sizes.
\end{definition}

\begin{definition}[Density profile of a sequence]
  \label{def:flag-density-seq}
  \lean{FlagAlgebras.flagDensitySeq}
  \uses{def:flag-seq, def:flag-density-n}
  \leanok
  $n \mapsto (F \mapsto d_1(F; s_n))$, the profile whose limit defines the
  semantics.
\end{definition}

\begin{definition}[Homomorphism as a PMF]
  \label{def:positivehom-topmf}
  \lean{FlagAlgebras.PositiveHom.toPMF}
  \uses{def:positive-hom, def:flag-with-size, def:basis-vector, lem:positivehom-basisvector-nonneg}
  \leanok
  The distribution on size-$\ell$ flags whose mass at $F$ is $\varphi[\![F]\!]$;
  $\varphi$'s random model at scale $\ell$.
\end{definition}

\begin{definition}[Flag-sequence measure]
  \label{def:flagseq-measure}
  \lean{FlagAlgebras.flagSeqMeasure}
  \uses{def:positive-hom, def:positivehom-topmf, def:flag-with-size}
  \leanok
  The independent product measure $\mu\{\varphi\}$ on sequences of flags, each
  component drawn from $\varphi$'s PMF; the source of the random flag sequence in
  Theorem~\ref{thm:hom-is-limit}.
\end{definition}

\begin{lemma}[Expected density is the downward ratio]
  \label{lem:integral-density-div}
  \lean{FlagAlgebras.integral_flagDensitySpace_eq_flagVectorDensity_div}
  \uses{def:finflag-pmf, def:flag-density-space, def:downward-factor, def:unlabel, def:label-extensions, thm:counting-identity, def:flag-density-n}
  \leanok
  The expected $F$-density under a large flag's measure equals the normalized
  downward density ratio; the finite-level identity behind Lemma~3.11.
\end{lemma}

\begin{theorem}[Lemma 3.11: convergence of expected densities]
  \label{thm:lemma-3-11}
  \lean{FlagAlgebras.tendsto_integral_flagDensitySpace_of_converge_flagSeq}
  \uses{def:converges-to, def:positive-hom-coe, lem:integral-density-div, def:downward}
  \leanok
  Along a flag sequence converging to $\varphi$, the expected $F$-densities tend
  to the normalized downward value
  $\varphi[\![\,[\![F]\!]\,]\!]_0 / \varphi[\![1]\!]_0$.
\end{theorem}

\begin{lemma}[Limit measure is supported on homomorphisms]
  \label{lem:measure-support}
  \lean{FlagAlgebras.flagSeq_limit_measure_support_positiveHomSpace}
  \uses{def:increases, def:positive-hom-space, thm:positive-hom-space-char}
  \leanok
  The weak limit of a convergent flag sequence's probability measures is
  concentrated on the positive-homomorphism space.
\end{lemma}

\begin{definition}[The measure $\mathbb{P}[\varphi_0]$]
  \label{def:prob-measure-witness}
  \lean{FlagAlgebras.probMeasure_extend_emptyType_positiveHom}
  \uses{thm:random-hom-measure, def:positive-hom-space}
  \leanok
  The chosen random-extension measure $\mathbb{P}[\varphi_0]$, the witness of
  Theorem~\ref{thm:random-hom-measure}.
\end{definition}

\begin{theorem}[Defining property of $\mathbb{P}[\varphi_0]$]
  \label{thm:prob-measure-spec}
  \lean{FlagAlgebras.probMeasure_extend_emptyType_positiveHom_spec}
  \uses{def:prob-measure-witness, thm:random-hom-measure, def:downward}
  \leanok
  Integrating $\varphi \mapsto \varphi(f)$ against $\mathbb{P}[\varphi_0]$
  recovers $\varphi_0[\![f]\!]_0 / \varphi_0[\![1]\!]_0$.
\end{theorem}

\begin{lemma}[Degenerate base homomorphism]
  \label{lem:downward-zero-at-hom}
  \lean{FlagAlgebras.downward_zero_at_hom}
  \uses{def:positive-hom, def:downward, def:basis-vector, thm:downward-linear}
  \leanok
  If $\varphi_0\langle\sigma\rangle_0 = 0$ then $\varphi_0[\![f]\!]_0 = 0$ for all
  $f$ --- the degenerate case handled separately in the cone and Cauchy--Schwarz
  arguments.
\end{lemma}

\begin{lemma}[Positive denominator]
  \label{lem:positivehom-one-downward-pos}
  \lean{FlagAlgebras.positiveHom_one_downward_pos}
  \uses{def:positive-hom, def:downward, lem:downward-factor-pos}
  \leanok
  If $\varphi_0\langle\sigma\rangle_0 > 0$ then $\varphi_0[\![1]\!]_0 > 0$, the
  denominator in the normalized downward value.
\end{lemma}


%%FILE blueprint/src/chapters/mantel.tex
\chapter{Application: Mantel's theorem}
\label{chap:mantel}

As a worked application of the whole framework, we prove \emph{Mantel's theorem}
--- the Tur\'an density of the triangle $K_3$ is $\tfrac12$. The flag-algebra
\emph{upper} bound (a sum-of-squares certificate, Chapter~\ref{chap:positivity})
and the extremal \emph{lower} bound meet in Theorem~\ref{thm:turan-k3}. (Source:
\texttt{LeanFlagAlgebras/MantelTheorem/}.)

\begin{definition}[Edge and triangle densities]
  \label{def:densities}
  \lean{MantelTheorem.K2, MantelTheorem.K3}
  \uses{def:fin-flag, def:basis-vector}
  \leanok
  $K_2, K_3 \in \mathcal{A}^\emptyset$ are the edge- and triangle-density basis
  elements. For a positive homomorphism $\varphi$, $\varphi(K_2)$ is the edge
  density and $\varphi(K_3)$ the triangle density.
\end{definition}

\begin{lemma}[Expansion of the edge density]
  \label{lem:expand-k2}
  \lean{MantelTheorem.expand_K2_on_three_vertex_graphs}
  \uses{def:densities, thm:density-chain-rules}
  \leanok
  Expanding onto three vertices,
  $K_2 = \tfrac13 E_3 + \tfrac23 P_3 + K_3$, where $E_3, P_3, K_3$ are the
  one-edge, path, and triangle densities.
\end{lemma}

\begin{lemma}[Square-positivity certificate]
  \label{lem:square-cert}
  \lean{MantelTheorem.FlagAlgebra_2_1_0_0_minus_FlagAlgebra_2_1_0_1_square_downward}
  \uses{def:downward, thm:square-downward-nonneg}
  \leanok
  Unlabelling the square of the difference of the two rooted two-vertex flags,
  \[ [\![\,(O_2^{1} - K_2^{1})^2\,]\!]_0
      = O_3 - \tfrac13 E_3 - \tfrac13 P_3 + K_3 \;\ge\; 0, \]
  a sum-of-squares certificate: the left side is a downward-of-square, hence
  nonnegative by Theorem~\ref{thm:square-downward-nonneg}.
\end{lemma}

\begin{theorem}[Mantel, flag-algebra form]
  \label{thm:mantel-flag}
  \lean{MantelTheorem.Mantel_theorem}
  \uses{def:densities}
  \leanok
  In $\mathcal{A}^\emptyset$, $K_2 \le \tfrac12\cdot 1 + K_3$.
\end{theorem}
\begin{proof}
  \uses{lem:expand-k2, lem:square-cert, lem:order-monotone}
  \leanok
  Substitute the expansion of Lemma~\ref{lem:expand-k2}, add the certificate of
  Lemma~\ref{lem:square-cert} and the nonnegativity of the three-vertex
  densities, and collect terms using the order arithmetic of
  Lemma~\ref{lem:order-monotone}.
\end{proof}

\begin{theorem}[Mantel, homomorphism form]
  \label{thm:mantel-hom}
  \lean{MantelTheorem.Mantel_theorem'}
  \uses{thm:mantel-flag, def:positive-hom}
  \leanok
  For every positive homomorphism $\varphi$ with $\varphi(K_3) = 0$ (a
  triangle-free limit), $\varphi(K_2) \le \tfrac12$.
\end{theorem}
\begin{proof}
  \uses{thm:mantel-flag}
  \leanok
  Apply $\varphi$ to Theorem~\ref{thm:mantel-flag}; positivity and
  $\varphi(K_3)=0$ give $\varphi(K_2) \le \tfrac12$.
\end{proof}

\begin{lemma}[Complete-bipartite lower bound]
  \label{lem:lower-bound}
  \lean{MantelTheorem.extremal_density_K3_ge}
  \leanok
  For $n \ge 2$, the maximal edge density of a triangle-free graph on $n$
  vertices is at least $\tfrac12$, witnessed by the balanced complete bipartite
  graph:
  $\operatorname{ex}(n,K_3)/\binom{n}{2} \ge \tfrac12$.
\end{lemma}

\begin{theorem}[Tur\'an density of the triangle]
  \label{thm:turan-k3}
  \lean{MantelTheorem.Turan_density_K3}
  \uses{thm:mantel-hom, lem:lower-bound}
  \leanok
  $\pi(K_3) = \tfrac12$.
\end{theorem}
\begin{proof}
  \uses{thm:mantel-hom, lem:lower-bound}
  \leanok
  The homomorphism form gives the upper bound $\pi(K_3) \le \tfrac12$; the
  complete-bipartite construction gives the matching lower bound. They meet at
  $\tfrac12$.
\end{proof}


%%FILE blueprint/src/chapters/overview.tex
\chapter{Overview}
\label{chap:overview}

This blueprint maps the core of the \texttt{LeanFlagAlgebras} library --- the
\texttt{FlagAlgebra/} directory --- which formalizes Razborov's flag algebras in
Lean~4. The development is layered, and each chapter below corresponds to one
layer:

\begin{itemize}
  \item \textbf{Flags} (Chapter~\ref{chap:flags}): labeled graphs, their
    isomorphisms, and the flag $\mathrm{Flag}\,\sigma$ as an isomorphism class.
  \item \textbf{Subflag densities} (Chapter~\ref{chap:density}): how often one
    flag occurs inside another, and the \emph{chain rules} relating single,
    pair, and triple densities.
  \item \textbf{The flag algebra $\mathcal{A}^\sigma$} (Chapter~\ref{chap:algebra}):
    the commutative $\mathbb{R}$-algebra obtained by quotienting formal flag
    combinations by Razborov's averaging relations, with multiplication defined
    through pair densities.
  \item \textbf{The unlabeling operator} (Chapter~\ref{chap:downward}):
    $[\![\,\cdot\,]\!]_0 : \mathcal{A}^\sigma \to \mathcal{A}^\emptyset$, which
    averages away the labels.
  \item \textbf{Positivity and SOS} (Chapter~\ref{chap:positivity}): positive
    homomorphisms, the semantic order $\le$ in which density bounds are stated,
    and the sum-of-squares certificates that prove them.
  \item \textbf{Limit semantics} (Chapter~\ref{chap:limits}): the analytic side
    (Razborov's Theorems 3.3 and 3.5) identifying positive homomorphisms with
    limits of subgraph densities, and the resulting Cauchy--Schwarz inequality.
  \item \textbf{Application} (Chapter~\ref{chap:mantel}): Mantel's theorem, the
    Tur\'an density of the triangle, proved with the machinery above.
\end{itemize}

Every boxed node in the dependency graph is a definition or theorem tagged with
the Lean declaration that realizes it; the graph edges are the mathematical
dependencies declared here. All declarations named live in the \texttt{FlagAlgebras}
namespace.


%%FILE blueprint/src/chapters/positivity.tex
\chapter{Positivity and sum-of-squares}
\label{chap:positivity}

Density bounds are inequalities in a semantic order on $\mathcal{A}^\sigma$
induced by its positive homomorphisms. Sum-of-squares certificates are the
mechanism that proves them. (Source: \texttt{FlagAlgebra/PositiveHom.lean},
\texttt{QuadraticForm.lean}.)

\begin{definition}[Positive homomorphism]
  \label{def:positive-hom}
  \lean{FlagAlgebras.PositiveHom}
  \uses{def:flag-algebra, def:hom}
  \leanok
  A \emph{positive homomorphism} is an $\mathbb{R}$-algebra homomorphism
  $\varphi : \mathcal{A}^\sigma \to \mathbb{R}$ that is nonnegative on every flag.
  These are the semantic evaluations of the algebra --- the limits of subgraph
  densities.
\end{definition}

\begin{definition}[Semantic cone]
  \label{def:semantic-cone}
  \lean{FlagAlgebras.semanticCone}
  \uses{def:positive-hom}
  \leanok
  The \emph{semantic cone} is $\{\, f \mid \varphi(f) \ge 0 \text{ for all
  positive } \varphi \,\}$. Membership is exactly an unconditional density
  inequality.
\end{definition}

\begin{definition}[Semantic order]
  \label{def:semantic-order}
  \lean{FlagAlgebras.le_def}
  \uses{def:semantic-cone}
  \leanok
  Define $f \le g$ to mean $g - f$ lies in the semantic cone, i.e.\ every
  positive homomorphism satisfies $\varphi(f) \le \varphi(g)$. This is a preorder
  on $\mathcal{A}^\sigma$, the order in which all bounds are stated.
\end{definition}

\begin{theorem}[Flags are nonnegative]
  \label{thm:flag-nonneg}
  \lean{FlagAlgebras.flag_geq_zero}
  \uses{def:semantic-order, def:basis-vector, lem:positivehom-basisvector-nonneg}
  \leanok
  Every flag satisfies $[\![F]\!] \ge 0$ --- the base positivity fact from which
  sum-of-squares bounds are assembled.
\end{theorem}

\begin{lemma}[Order arithmetic]
  \label{lem:order-monotone}
  \lean{FlagAlgebras.flag_add_le_add, FlagAlgebras.nonneg_smul_nonneg_geq_zero}
  \uses{def:semantic-order}
  \leanok
  The semantic order is compatible with the algebra operations: inequalities may
  be added termwise, and a nonnegative scalar times a nonnegative element is
  nonnegative.
\end{lemma}

\begin{definition}[Flag quadratic form]
  \label{def:flag-quadratic-form}
  \lean{FlagAlgebras.flagQuadraticForm}
  \uses{def:flag-algebra, def:flag-algebra-vec}
  \leanok
  Given a matrix $M$ and a vector $v = (v_1,\dots,v_n)$ of flag-algebra elements,
  \[ \mathrm{flagQuadraticForm}\,M\,v \;=\; \sum_{i,j} M_{ij}\,(v_i\, v_j), \]
  the raw sum-of-squares certificate object.
\end{definition}

\begin{theorem}[PSD forms are nonnegative]
  \label{thm:sos-nonneg}
  \lean{FlagAlgebras.flagQuadraticForm_nonneg}
  \uses{def:flag-quadratic-form, def:semantic-cone, def:semantic-order}
  \leanok
  If $M \succeq 0$ (positive semidefinite) then
  $\mathrm{flagQuadraticForm}\,M\,v \ge 0$ in the semantic cone.
\end{theorem}
\begin{proof}\leanok
  Apply any positive homomorphism $\varphi$: it turns the form into
  $\sum_{ij} M_{ij}\,\varphi(v_i)\varphi(v_j) = x^\top M x \ge 0$ with
  $x_i = \varphi(v_i)$, since $M \succeq 0$.
\end{proof}

\begin{theorem}[Unlabeled PSD forms are nonnegative]
  \label{thm:sos-downward-nonneg}
  \lean{FlagAlgebras.flagQuadraticForm_downward_nonneg}
  \uses{def:flag-quadratic-form, thm:sos-nonneg, def:downward}
  \leanok
  The unlabeled form $[\![\,\mathrm{flagQuadraticForm}\,M\,v\,]\!]_0 \ge 0$ for
  $M \succeq 0$. This is the certificate actually used, since density bounds live
  in the unlabeled algebra $\mathcal{A}^\emptyset$.
\end{theorem}
\begin{proof}\leanok
  \uses{thm:downward-preserve-cone}
  The form is in the semantic cone by Theorem~\ref{thm:sos-nonneg}, and downward
  preserves the semantic cone (Theorem~\ref{thm:downward-preserve-cone}), so its
  unlabeling is nonnegative.
\end{proof}

\section{Supporting constructions}

\begin{definition}[Algebra homomorphism type]
  \label{def:hom}
  \lean{FlagAlgebras.Hom}
  \uses{def:flag-algebra}
  \leanok
  $\mathrm{Hom}\,\sigma = \mathcal{A}^\sigma \to_{\mathbb{R}\text{-alg}}
  \mathbb{R}$, the underlying $\mathbb{R}$-algebra homomorphisms; a positive
  homomorphism is such a map that is additionally nonnegative on flags.
\end{definition}

\begin{definition}[Vector of algebra elements]
  \label{def:flag-algebra-vec}
  \lean{FlagAlgebras.FlagAlgebraVec}
  \uses{def:flag-algebra}
  \leanok
  $\mathrm{FlagAlgebraVec}\,\sigma\,n = \mathrm{Fin}\,n \to \mathcal{A}^\sigma$,
  the variables of a quadratic form.
\end{definition}

\begin{lemma}[Homomorphisms are nonnegative on flags]
  \label{lem:positivehom-basisvector-nonneg}
  \lean{FlagAlgebras.positiveHom_basisVector_ge_zero}
  \uses{def:positive-hom, def:basis-vector}
  \leanok
  $0 \le \varphi[\![F]\!]$ for every positive homomorphism $\varphi$ and flag
  $F$ --- the defining positivity, restated as $0 \le \cdot$.
\end{lemma}

\begin{lemma}[Order via nonnegative difference]
  \label{lem:flag-sub-nonneg}
  \lean{FlagAlgebras.flag_sub_nonneg}
  \uses{def:semantic-order}
  \leanok
  $f \le g \iff 0 \le g - f$.
\end{lemma}

\begin{lemma}[Translation invariance of the order]
  \label{lem:order-monotone-sides}
  \lean{FlagAlgebras.flag_add_le_add_left, FlagAlgebras.flag_add_le_add_right}
  \uses{lem:order-monotone}
  \leanok
  Adding a fixed element on either side preserves the semantic order; these
  register as the \texttt{AddLeftMono}/\texttt{AddRightMono} instances.
\end{lemma}


%%FILE blueprint/src/content.tex
% Master content file for the Lean Flag Algebras blueprint.
% Each \input is one chapter = one layer of the FlagAlgebra/ core library,
% in dependency order, followed by Mantel's theorem as a worked application.

\input{chapters/overview}
\input{chapters/flags}
\input{chapters/density}
\input{chapters/algebra}
\input{chapters/downward}
\input{chapters/positivity}
\input{chapters/limits}
\input{chapters/compute}
\input{chapters/mantel}

